\chapter{Cinematic functions and their maximal estimate}
\label{chap:cinematic}

This chapter follows Pramanik, Yang, and Zahl~\cite{PYZ}. The
application to sticky Kakeya uses its maximal estimate for a
bounded family of functions produced from the nonsingular slope.
The metric on $C^2(I)$ is the maximum of the sup norms of a
function and its first two derivatives.

\section{The maximal estimate}

\begin{definition}[One-dimensional sets and quasi-products; PYZ, Definitions 1.4--1.5]
\label{def:quasi-product}
\lean{Kakeya.Cinematic.IsDeltaSet,Kakeya.Cinematic.IsQuasiProduct}
\leanok
A $(\delta,\alpha;C)_1$-set $A\subset[0,1]$ is a union of
intervals of length $\delta$ such that
\[
                         |A\cap I|\leq C\delta^{1-\alpha}|I|^\alpha
\]
for every interval $I\subset[0,1]$. A Borel set $E\subset[0,1]^2$
is a $(\delta,\alpha;C)_1\times(\delta,\beta;C)_1$ quasi-product
if $E=\bigcup_{a\in A}\{a\}\times B_a$, where $A$ and each
$B_a$ have the respective one-dimensional conditions.
\end{definition}

\begin{definition}[Cinematic families; PYZ, Definition 1.6]
\label{def:cinematic-family}
\lean{Kakeya.Cinematic.IsCinematicFamily}
\leanok
A family $\cF\subset C^2(I)$ is cinematic with constants $K,D$
if it has diameter at most $K$, is a doubling metric space with
doubling constant at most $D$, and for all $f,g\in\cF$,
\[
 \inf_{t\in I}\bigl(|f(t)-g(t)|+|f'(t)-g'(t)|+|f''(t)-g''(t)|\bigr)
                                      \geq K^{-1}\|f-g\|_{C^2}.
\]
\end{definition}

\begin{theorem}[Cinematic maximal estimate; PYZ, Theorem 1.7, normalized formulation]
\label{thm:pyz}
\lean{Kakeya.Cinematic.Theorem1_7UniformC2Statement}
\uses{def:cinematic-family,def:quasi-product}
Fix $D,K\geq1$ and an absolute bound
$\sup_{f\in\cF}\|f\|_{C^2}\leq M$. For every $\eps>0$,
all sufficiently small $\delta>0$ have the following property.
Let $0<\alpha\leq\zeta\leq1$, let $E$ be a
$(\delta,\alpha;\delta^{-\eps})_1\times
(\delta,\alpha;\delta^{-\eps})_1$ quasi-product, and let
$F\subset\cF$ be finite with
\[
           \#(F\cap B(g,r))\leq\delta^{-\eps}(r/\delta)^\zeta
           \quad(r\geq\delta).
\]
Writing $f^\delta=N_\delta(\graph f)$, one has
\[
       \int_E\left(\sum_{f\in F}\one_{f^\delta}\right)^{3/2}
             \leq\delta^{2-\alpha/2-\zeta/2-C\eps}\#F,
\]
where $C$ depends only on $D$ and the small-scale threshold may
also depend on $K,M,\eps$.
\end{theorem}
\begin{proof}
\uses{lem:pyz-local-maximal,prop:bipartite-tangency}
Reduce to a short interval, select a multiplicity level, and
pigeonhole pairwise distances and tangency parameters. The
two-ends selection prevents most functions from lying in a much
smaller function-space ball. Assign common tangent rectangles to
the remaining points. The bipartite rectangle estimate bounds
their number, while the quasi-product condition bounds the mass
inside each rectangle. Summing the multiplicity levels gives the
$3/2$-moment estimate, and a finite interval covering gives the
global result.
\end{proof}

\begin{remark}[Uniformity in the family]
\label{rem:pyz-uniformity}
The printed theorem says that the threshold depends only on
$D,K,\eps$. The cinematic conditions control differences between
functions, not their absolute $C^2$ norms. A singleton containing
an arbitrarily rapidly oscillating function still has diameter
zero and satisfies the pairwise condition. At a prescribed scale
its Euclidean graph neighbourhood can cover a fixed part of the
unit square. Thus a threshold uniform over all such singletons
needs an additional absolute bound, or must depend on the family.
The Lean formulation fixes $M$ before choosing $\delta_0$.

The general normalized statement is retained as a proposition in
the repository. The closed estimate consumed by sticky Kakeya is
the $L^{3/2}$ interface of Theorem~\ref{thm:cinematic-l32}, proved
in \texttt{wz2\_uniformC2\_input}; its fixed constants and
quantifiers are recorded there. This distinction is why the full
$\alpha,\zeta$ statement above has no proof-completion mark.
\end{remark}

\begin{proposition}[The curve-generated family; PYZ, Proposition 2.1]
\label{prop:pyz-curve-family}
\uses{def:quasi-product,def:cinematic-family}
Let $\gamma:I\to S^2$ be $C^2$ with
$\det(\gamma,\gamma',\gamma'')\ne0$ everywhere. Fix
$0<\alpha\leq\zeta\leq1$ and $\eps>0$.
For a $\delta$-separated set $Z\subset B(0,1)\subset\R^3$
with $\#(Z\cap B(z,r))\leq\delta^{-\eps}(r/\delta)^\zeta$,
put $\Gamma_z=\graph(t\mapsto z\cdot\gamma(t))$.
At sufficiently small $\delta$, a quasi-product $E$ as in
Theorem~\ref{thm:pyz} satisfies
\[
 \int_E\left(\sum_{z\in Z}\one_{N_\delta(\Gamma_z)}\right)^{3/2}
              \leq\delta^{2-\alpha/2-\zeta/2-C\eps}\#Z.
\]
\end{proposition}
\begin{proof}
\uses{thm:pyz}
The determinant condition makes the two-jet of
$z\cdot\gamma$ quantitatively determine $z$. Thus the map into
$C^2(I)$ is bi-Lipschitz on the bounded parameter set, with a
uniform absolute $C^2$ bound supplied by $\gamma$. Apply the
cinematic estimate after splitting into a fixed number of patches.
\end{proof}

\section{Geometry on a short interval}

Fix a sufficiently short interval $J\subset I$, with its length
controlled by $K$. In the statements below, $aJ$ is the interval
with the same midpoint and $a$ times the length. Constants in
geometric comparisons may depend on $K$.

\begin{definition}[Curvilinear rectangles and tangency; PYZ, Definitions 3.1--3.4]
\label{def:cinematic-rectangle}
\lean{Kakeya.Cinematic.CurvilinearRectangle,Kakeya.Cinematic.CurvilinearRectangle.IsLambdaTangent,Kakeya.Cinematic.CurvilinearRectangle.AreLambdaComparable}
\leanok
For $0<\delta\leq t\leq1$, a $(\delta,t)$-rectangle is
\[
         f^\delta(L)=\{(\theta,y):\theta\in L, |y-f(\theta)|\leq\delta\},
         \qquad |L|=\sqrt{\delta/t},
\]
with $f\in\cF$. Two rectangles are $\lambda$-comparable if
their union lies in a $(\lambda\delta,t)$-rectangle. A rectangle
is $\lambda$-tangent to $g$ if it lies in the vertical
$\lambda\delta$-neighbourhood of $g$; unqualified tangency means
$\lambda=5$. Two function sets $W,B$ are $t$-separated if
$\|w-b\|_{C^2}\geq t$ for all $w\in W,b\in B$.
\end{definition}

\begin{lemma}[The short-interval dichotomy; PYZ, Lemma 3.5]
\label{lem:cinematic-dichotomy}
\lean{Kakeya.Cinematic.PreliminaryDichotomyStatement}
\uses{def:cinematic-family}
Let $L\subset J$ have length at most $(6K)^{-1}$. For
$k=f-g$ or $k=f'-g'$, either $|k|<(3K)^{-1}\|f-g\|_{C^2}$
everywhere on $L$, or
$|k|\geq(6K)^{-1}\|f-g\|_{C^2}$ everywhere there.
If the small alternative holds for both differences, the large
alternative holds for $f''-g''$.
\end{lemma}
\begin{proof}
The derivative of each of the first two differences is bounded by
$\|f-g\|_{C^2}$. Hence it cannot vary by more than
$(6K)^{-1}\|f-g\|_{C^2}$ on $L$. The final assertion follows
from cinematic nondegeneracy.
\end{proof}

\begin{lemma}[At most two intersections; PYZ, Lemma 3.6]
\label{lem:two-zeros}
\lean{Kakeya.Cinematic.cinematic_difference_has_at_most_two_zeros}
\leanok
\uses{def:cinematic-family}
For distinct $f,g\in\cF$, the difference $f-g$ has at most
two zeros on $J$.
\end{lemma}
\begin{proof}
\uses{lem:cinematic-dichotomy}
\leanok
Three zeros of $f-g$ would give two zeros of $f'-g'$ and a
zero of $f''-g''$ by Rolle's theorem. The dichotomy and
nondegeneracy exclude this.
\end{proof}

\begin{lemma}[Two derivative zeros; PYZ, Lemma 3.6, second assertion]
\label{lem:two-derivative-zeros}
\uses{def:cinematic-family}
If $f'-g'$ has two distinct zeros on $J$, then
$|f-g|\geq(6K)^{-1}\|f-g\|_{C^2}$ throughout $J$.
\end{lemma}
\begin{proof}
\uses{lem:cinematic-dichotomy}
Rolle's theorem gives a zero of $f''-g''$. The dichotomy forces
$f'-g'$ into its small alternative. Nondegeneracy then excludes
the small alternative for $f-g$. This additional assertion in the
paper is not part of the linked Lean \texttt{TwoZerosStatement}.
\end{proof}

\begin{definition}[Tangency parameter; PYZ, Definition 3.7]
\label{def:tangency-parameter}
\lean{Kakeya.Cinematic.tangencyParameterOn}
\leanok
Set
\[
          \Delta(f,g)=\min_{\theta\in J/2}
                        (|f-g|(\theta)+|f'-g'|(\theta)).
\]
The minimum exists by continuity on a compact interval.
\end{definition}

\begin{lemma}[Tangency and sublevel intervals; PYZ, Lemma 3.8]
\label{lem:tangency-sublevels}
\lean{Kakeya.Cinematic.TangencySublevelDiameterStatement,Kakeya.Cinematic.tangency_sublevel_diameter}
\uses{def:tangency-parameter}
Write $h=f-g$, $t=\|f-g\|_{C^2}>0$, and $\Delta=\Delta(f,g)$.
If $\Delta\leq c_1(K)t$, there is a unique critical point
$\theta_0\in3J/5$ with $|h(\theta_0)|\lesssim\Delta$.
Outside $4J/5$ one has $|h|\gtrsim t$; if $h$ changes sign,
its two zeros lie in $4J/5$ at distance
$O(\sqrt{\Delta/t})$ from $\theta_0$.
For $0<\delta\leq t/(6K)$, the set
$\{\theta\in J/4:|h(\theta)|\leq\delta\}$ lies in an
interval of length $O(\sqrt{(\Delta+\delta)/t})$ and has at
most two interval components, each of length
\[
                         O\left(\frac{\delta}{\sqrt{(\Delta+\delta)t}}\right).
\]
A component containing a point where $|h|\leq\delta/2$ has
the corresponding lower length bound, subject to the interior
window in the construction.
\end{lemma}
\begin{proof}
\uses{lem:cinematic-dichotomy,lem:two-zeros}
The dichotomy reduces the problem to either a monotone difference
with controlled first derivative or a strictly convex or concave
difference with second derivative comparable to $t$.
The linear and quadratic estimates in these cases give the
location, diameter, and component-length bounds.
\end{proof}

\begin{lemma}[A common tangent rectangle; PYZ, Lemma 3.9]
\label{lem:common-tangent}
\lean{Kakeya.Cinematic.CommonTangentRectangleRobustStatement,Kakeya.Cinematic.common_tangent_rectangle_robust}
\uses{def:cinematic-rectangle,def:tangency-parameter}
If a $(\delta,t)$-rectangle over $J/4$ is tangent to both
$f$ and $g$, then
\[
                         (\Delta(f,g)+\delta)\|f-g\|_{C^2}
                                           \lesssim\delta t.
\]
For a general tangency dilation, the implicit constant is
polynomial in that dilation.
\end{lemma}
\begin{proof}
\uses{lem:tangency-sublevels}
The common rectangle forces the sublevel set of $f-g$ to have
length comparable to $\sqrt{\delta/t}$. Compare this with
the component-length upper bound and rearrange.
\end{proof}

\begin{lemma}[Enlarging a small-value interval; PYZ, Lemma 3.10]
\label{lem:cinematic-interval-growth}
\uses{def:cinematic-family}
If $L\subset J/4$, $0<\delta\leq\|f-g\|_{C^2}/(6K)$,
and $|f-g|\leq\delta$ on $L$, then for $\lambda\geq1$,
$|f-g|\lesssim\lambda^2\delta$ on $\lambda L\cap J$.
\end{lemma}
\begin{proof}
\uses{lem:tangency-sublevels}
Use the linear or quadratic description of the difference near
its small-value interval. Enlarging the horizontal interval
by $\lambda$ increases the controlling quadratic by at most
a factor comparable to $\lambda^2$.
\end{proof}

\begin{definition}[Dilation of a rectangle; PYZ, Definition 3.11]
\label{def:rectangle-dilation}
For $R=f^\delta(L)$ define
$\lambda R=f^{\lambda\delta}(\lambda L\cap J)$.
This dilation enlarges the vertical and horizontal parameters
separately; it is not Euclidean homothety about a point.
\end{definition}

\begin{lemma}[Comparable rectangles; PYZ, Lemma 3.12]
\label{lem:comparable-rectangles}
\lean{Kakeya.Cinematic.ComparableRectanglesStatement,Kakeya.Cinematic.comparable_rectangles_have_close_graphs}
\uses{def:cinematic-rectangle}
If $f^\delta(L)$ and $g^\delta(L')$ are
$\lambda$-comparable $(\delta,t)$-rectangles over $J/4$,
the convex hull $L''$ of $L\cup L'$ has length at most
$\sqrt{\lambda\delta/t}$ and
$|f-g|\lesssim\lambda^3\delta$ on $L''$.
\end{lemma}
\begin{proof}
\uses{lem:cinematic-interval-growth}
The containing rectangle bounds the horizontal hull. Compare
each graph to its central graph on its own interval, then
extend the comparison across the hull using the interval-growth
estimate.
\end{proof}

\begin{corollary}[Transitivity up to a constant; PYZ, Corollary 3.13]
\label{cor:rectangle-transitivity}
\lean{Kakeya.Cinematic.ComparabilityTransitivityStatement}
\uses{def:cinematic-rectangle}
If $R_1$ is $\lambda_1$-comparable to $R_2$ and $R_2$
is $\lambda_2$-comparable to $R_3$, then $R_1$ and $R_3$
are $\lambda_3$-comparable for a constant depending only
on $K,\lambda_1,\lambda_2$.
\end{corollary}
\begin{proof}
\uses{lem:comparable-rectangles}
Combine the graph and interval-hull bounds for the two pairs
and choose a containing rectangle with the resulting constant.
\end{proof}

\begin{lemma}[Shrinking incomparable rectangles; PYZ, Lemma 3.14]
\label{lem:rectangle-shrink}
\uses{def:cinematic-rectangle}
For $0<c<1$, if the original rectangles are
$C(c,K)$-incomparable, then shrinking their horizontal
intervals by $c$ leaves them $100$-incomparable.
\end{lemma}
\begin{proof}
\uses{lem:cinematic-interval-growth}
Otherwise enlarge the common containing rectangle for the
shrunk pair. The graph-growth bound would make the original
pair comparable with a constant depending only on $c,K$.
\end{proof}

\begin{lemma}[Packing inside one rectangle; PYZ, Lemma 3.15]
\label{lem:rectangle-packing}
\lean{Kakeya.Cinematic.RectanglePackingStatement,Kakeya.Cinematic.incomparable_rectangles_packing_bound}
\leanok
\uses{def:cinematic-rectangle}
Let the central functions of a pairwise $100$-incomparable
family of $(\delta,t)$-rectangles lie in a $C^2$ ball of
radius $3t$. If all its rectangles lie in one
$(\lambda\delta,t)$-rectangle, $\lambda\geq100$, then
their number is $O_K(\lambda^{5/2})$.
\end{lemma}
\begin{proof}
\uses{lem:comparable-rectangles}
\leanok
Discretize the possible interval centres and the graph values
and slopes inside the containing rectangle. Two rectangles
with the same discretized data would be comparable. Counting
the admissible data gives the stated power of $\lambda$.
\end{proof}

\begin{lemma}[Increasing incomparability; PYZ, Lemma 3.16]
\label{lem:increase-incomparability}
\uses{def:cinematic-rectangle}
Under the same function-ball hypothesis, a finite pairwise
$100$-incomparable family has, for any fixed $C\geq100$,
a pairwise $C$-incomparable subfamily of comparable size.
The comparison constant may depend on $C,K$.
\end{lemma}
\begin{proof}
\uses{lem:rectangle-packing}
Choose a maximal $C$-incomparable subfamily. The packing
bound limits how many original rectangles can be assigned
to each selected rectangle.
\end{proof}

\begin{lemma}[A coarse count of rectangles; PYZ, Lemma 3.17]
\label{lem:rectangle-coarse-count}
\uses{def:cinematic-rectangle,def:cinematic-family}
For $\lambda\geq100$, a pairwise $\lambda$-incomparable
family of $(\delta,t)$-rectangles in a cinematic family
has at most $\delta^{-C(D)}$ elements, at sufficiently
small $\delta$ depending on the fixed geometric constants.
\end{lemma}
\begin{proof}
\uses{lem:rectangle-packing}
Use the doubling property to cover function space by a
polynomial number of small balls. Within each ball, count
the possible interval locations and apply the packing bound.
\end{proof}

\section{Bipartite tangencies and lenses}

\begin{proposition}[Bipartite tangency bound; PYZ, Proposition 4.1]
\label{prop:bipartite-tangency}
\lean{Kakeya.Cinematic.BipartiteTangencyStatement,Kakeya.Cinematic.bipartite_tangency_robust_full}
\uses{def:cinematic-rectangle}
Let $A\geq1$, $0<\delta\leq At$, and let finite function
families $W,B$ have union of diameter at most $6t$ and mutual
separation at least $t/A$. Let $\mathcal R$ be a pairwise
$100$-incomparable collection of $(\delta,t)$-rectangles over
$J/4$, each tangent to at least $\mu\geq1$ members of $W$
and $\nu\geq1$ members of $B$. Then
\[
 \#\mathcal R\leq CA^C
       \left(\frac{\#W}{\mu}+\frac{\#B}{\nu}\right)^{3/2}
       \log\left(\frac{\#W}{\mu}+\frac{\#B}{\nu}\right),
\]
where $C=C(K)$. The formal rectangle interface keeps the
admissible scale and dilation conditions explicit.
\end{proposition}
\begin{proof}
\uses{lem:bipartite-perturbation,lem:large-delta-reduction,
  lem:separation-reduction,lem:incomparability-reduction}
First reduce to well-separated colour classes, a small ratio
$\delta/t$, and sufficiently strong incomparability. For
$\mu=\nu=1$, associate rectangles to nonoverlapping lenses
after a small perturbation and apply the lens count. For
larger multiplicities, independently sample the two colour
classes at rates comparable to $1/\mu$ and $1/\nu$ and
average the resulting unit-multiplicity bound.
\end{proof}

\begin{lemma}[Reduction to small $\delta/t$; PYZ, Lemma 4.2]
\label{lem:large-delta-reduction}
If the bipartite estimate holds for $\delta\leq ct$, for
a fixed $c>0$, then it holds for $\delta\leq At$ after
changing its constant by a quantity depending on $K,c$ and
the original constant.
\end{lemma}
\begin{proof}
\uses{lem:rectangle-packing}
In the remaining range the rectangles have a controlled
horizontal length. A packing argument gives a polynomial
bound in the scale ratio, which is absorbed by the factor
$A^C$.
\end{proof}

\begin{lemma}[Reduction to separated colour classes; PYZ, Lemma 4.3]
\label{lem:separation-reduction}
If the bipartite estimate holds for $2t$-separated pairs,
then it holds for $t/A$-separated pairs of union diameter
at most $6t$, with a polynomial loss in $A$.
\end{lemma}
\begin{proof}
Partition function space into smaller balls and pair those
which meet the two colour classes. Rescale each pair to the
well-separated configuration and sum the estimates. The
geometric tangency constant must also be rescaled.
\end{proof}

\begin{lemma}[Reduction of incomparability; PYZ, Lemma 4.4]
\label{lem:incomparability-reduction}
It suffices to prove the bipartite estimate for pairwise
$2C_c$-incomparable rectangles, where $C_c=C_c(K)$ is the
constant in the overlapping-lens comparison below.
\end{lemma}
\begin{proof}
\uses{lem:increase-incomparability}
Select a strongly incomparable subfamily of comparable size
and apply the estimate to it.
\end{proof}

\begin{definition}[Pseudo-circles and lenses; PYZ, Definitions 4.5--4.7]
\label{def:pseudocircles}
A family of closed Jordan curves consists of pseudo-circles
if two distinct curves meet at most twice and all their
intersections are proper crossings. Two crossings cut the
curves into arcs. A lens is a simple closed curve formed by
one arc from each curve. Lenses overlap when they share an
arc of positive length, not merely an endpoint.
\end{definition}

\begin{theorem}[Lens count; PYZ, Theorem 4.8]
\label{thm:lens-count}
\lean{Kakeya.Cinematic.graph_lens_bound_from_sequences}
\uses{def:pseudocircles}
A family of $n$ pseudo-circles supports at most
$O(n^{3/2}\log(2+n))$ pairwise nonoverlapping lenses.
The formalized application uses the equivalent graph-lens
count on a short interval.
\end{theorem}
\begin{proof}
Order the intersection arcs along the curves and encode
the nonoverlapping lenses by the resulting finite incidence
sequences. The combinatorial lens estimate bounds their
number. Closing the graph arcs outside the working interval
reduces the geometric formulation to this count.
\end{proof}

\begin{lemma}[Overlapping lenses force comparable rectangles; PYZ, Lemma 4.9]
\label{lem:lens-comparability}
\lean{Kakeya.Cinematic.common_tangent_overlap_comparable}
\uses{def:pseudocircles,def:cinematic-rectangle}
Suppose $\delta/t$ is sufficiently small. Let $w,b_1,b_2$
belong to the cinematic family, with $\|w-b_i\|_{C^2}\geq t$,
and suppose $w,b_i$ form a lens on $J$. If a
$(\delta,t)$-rectangle $R_i$ is $2\lambda$-tangent to
$w,b_i$ for each $i$, and the two lenses overlap, then
$R_1,R_2$ are $C_c(K,\lambda)$-comparable.
\end{lemma}
\begin{proof}
\uses{lem:tangency-sublevels,lem:common-tangent}
The common lens arc forces the two tangency windows to
occur near the same critical region of the graph
differences. The sublevel and common-tangent estimates
bound both their separation and their vertical discrepancy.
\end{proof}

\begin{lemma}[The generic unit-multiplicity case; PYZ, Lemma 4.10]
\label{lem:generic-bipartite}
\uses{def:cinematic-rectangle}
Let $W,B$ be $t$-separated, $\delta\leq c_2t$, and let
$\mathcal R$ be pairwise $C_c$-incomparable rectangles,
each assigned a tangent pair $(w_R,b_R)$. Assume every
assigned pair crosses twice in $J$, all functions have
distinct values at each endpoint of $J$, and no two
functions have an exact tangency in $J$. Then
\[
                \#\mathcal R\lesssim_K
                (\#W+\#B)^{3/2}\log(2+\#W+\#B).
\]
\end{lemma}
\begin{proof}
\uses{thm:lens-count,lem:lens-comparability}
Complete the graph arcs to pseudo-circles. The assigned
pairs give lenses, and incomparability prevents them from
overlapping. Apply the lens count.
\end{proof}

\begin{lemma}[Removing genericity; PYZ, Lemma 4.11]
\label{lem:bipartite-perturbation}
\uses{def:cinematic-rectangle}
The same unit-multiplicity bound holds without the
genericity assumptions if $W,B$ are $2t$-separated,
$\delta\leq c_2t/2$, and the rectangles are pairwise
$2C_c$-incomparable and tangent to both colour classes.
\end{lemma}
\begin{proof}
\uses{lem:generic-bipartite,cor:common-shift}
Choose a common vertical shift retaining a fixed fraction
of the assigned rectangles. Perturb the finite family
slightly to separate endpoint values and remove exact
tangencies. The margins in the hypotheses preserve the
required incidences and incomparability. Apply the generic
bound and pass to the original configuration.
\end{proof}

\begin{lemma}[Creating two crossings by a shift; PYZ, Lemma 4.12]
\label{lem:two-crossing-shift}
\uses{def:cinematic-rectangle}
If $\|f-g\|_{C^2}\geq2t$, $0<\delta<c_1t$, and
$f,g$ are tangent to a common $(\delta,t)$-rectangle,
then for some $e\in\{-1,0,1\}$ the function
$f+e\lambda\delta$ crosses $g$ twice in $J$ and is
$2\lambda$-tangent to that rectangle, where
$\lambda=\lambda(K)$ is fixed sufficiently large.
\end{lemma}
\begin{proof}
\uses{lem:tangency-sublevels}
The common tangency forces a small extremum of $f-g$.
One of the three shifts puts this extremum on the opposite
side of zero from the outer part of the interval. Strict
curvature gives exactly two crossings.
\end{proof}

\begin{corollary}[One shift for many rectangles; PYZ, Corollary 4.13]
\label{cor:common-shift}
\uses{def:cinematic-rectangle}
In the setting of Lemma~\ref{lem:bipartite-perturbation},
one common shift $e\lambda\delta$ of the first colour
class retains at least one third of the rectangles, with
each retained rectangle tangent to an assigned pair
crossing twice in $J$.
\end{corollary}
\begin{proof}
\uses{lem:two-crossing-shift}
Assign a successful shift to each rectangle and choose
the largest of the three classes.
\end{proof}

\section{The local maximal argument}

\begin{lemma}[The short-curve estimate; PYZ, Lemma 5.1]
\label{lem:pyz-local-maximal}
\uses{def:cinematic-family,def:quasi-product}
Under the normalized hypotheses of Theorem~\ref{thm:pyz},
let $J'=J/16$ and $E'=E\cap(J'\times\R)$. Then
\[
       \int_{E'}\left(\sum_{f\in F}\one_{f^\delta}\right)^{3/2}
           \leq\delta^{2-\alpha/2-\zeta/2-C_0\eps}\#F,
\]
with $C_0=C_0(D)$ and the small-scale threshold depending
on the fixed absolute $C^2$ bound as well.
\end{lemma}
\begin{proof}
\uses{lem:cinematic-two-ends,lem:fine-rectangle-assignment,
  lem:fine-rectangle-selection,lem:good-function-pairs,
  lem:pair-rectangle-count,lem:coarse-fibre-nonconcentration,
  prop:bipartite-tangency}
Select a level set of the graph multiplicity. A two-ends
selection fixes a function-space distance scale $t$ and a
tangency scale $\Delta$. Cover the retained points by
essentially disjoint fine rectangles, and group those into
coarse rectangles. The quasi-product bound controls how
many fine rectangles can lie inside a coarse one by their
mass. Counting well-separated pairs of tangent functions
gives a second bound. Interpolating the two bounds with
weights $1/4$ and $3/4$ produces the exponent $3/2$.
Finally use the bipartite estimate for the coarse
rectangles and sum the multiplicity levels.
\end{proof}

We record the numbered selections used in that proof. At
this stage $F_B$ lies in a function-space ball of radius
$3t$. The set $F_B(x)$ consists of the functions retained
at a point $x$ by the distance selection; $Y_1(f)$ is the
refined point-function relation after tangency selection.
The retained point set is $E_2$, and $\Delta$ is the
common tangency scale. For each $x\in E_2$, write $R_x$
for its fine rectangle. The sets $Y_2(R)$ are the
enlarged rectangle shadings formed from this relation.
A good pair $(f,R)$ means $x(R)\in Y_1(f)$, where
$x(R)$ is its assigned point. Subsequent primes indicate
the common degree refinements in PYZ Section 5.

\begin{lemma}[Nonconcentration after two-ends selection; PYZ, Lemma 5.2]
\label{lem:cinematic-two-ends}
\lean{Kakeya.Cinematic.TwoEndsSelectionStatement}
For each retained point $x$ and every $g\in C^2(J)$,
if $\delta/t<\lambda<1$, then
\[
 \#\{f\in F_B(x):\|f-g\|_{C^2}\leq\lambda t\}
               \leq4(2\lambda)^\eps\#F_B(x).
\]
\end{lemma}
\begin{proof}
Choose the distance scale where the normalized local
count is largest. Excess concentration at a smaller
radius would contradict that choice.
\end{proof}

\begin{lemma}[Assigning a fine rectangle; PYZ, Lemma 5.3]
\label{lem:fine-rectangle-assignment}
\lean{Kakeya.Cinematic.FineRectangleAssignmentStatement}
For $t'=C_Rt\Delta/\delta$, with $C_R=C_R(K)$
sufficiently large, each $x\in E_2$ lies in a
$(\delta,t')$-rectangle over $J/8$ to which every
$f$ with $x\in Y_1(f)$ is tangent.
\end{lemma}
\begin{proof}
\uses{lem:tangency-sublevels}
The common distance and tangency bounds give a common
sublevel interval of the required length. Choose the
rectangle around a retained reference function through
$x$ and use the sublevel estimates for the others.
\end{proof}

\begin{lemma}[Compatibility of rectangle shadings; PYZ, Lemma 5.4]
\label{lem:rectangle-shading}
For a sufficiently large enlargement constant in the
definition of $Y_2$, $100$-comparable rectangles satisfy
$Y'_2(R')\subset Y_2(R)$.
\end{lemma}
\begin{proof}
\uses{lem:comparable-rectangles}
Comparability controls the interval hull and the
difference of the central graphs. The chosen enlargement
contains the corresponding smaller shading.
\end{proof}

\begin{lemma}[Selecting disjoint fine rectangles; PYZ, Lemma 5.5]
\label{lem:fine-rectangle-selection}
\lean{Kakeya.Cinematic.maximal_fine_rectangle_selection}
There is a pairwise $100$-incomparable subcollection
$\mathcal R$ of the rectangles $R_x$ and a number
$\Lambda>0$ such that
\[
 \Lambda\leq|Y_2(R)|<2\Lambda,\qquad
 |E_2|\lessapprox\sum_{R\in\mathcal R}|Y_2(R)|
                          \asymp\Lambda\#\mathcal R.
\]
The retained shadings are disjoint, and
$\#\mathcal R\leq\delta^{-C}$,
$\Lambda\geq\delta^{C'}$ for constants depending on
the fixed family parameters.
\end{lemma}
\begin{proof}
\uses{lem:rectangle-shading,lem:rectangle-coarse-count}
Greedily select incomparable rectangles and assign the
remaining points to them. Pigeonhole the shading
measures and disjointify while retaining a fixed
fraction. The coarse rectangle count bounds the number
of dyadic measure classes that matter.
\end{proof}

\begin{lemma}[Tangency to the coarse rectangle; PYZ, Lemma 5.6]
\label{lem:coarse-tangency}
If $R=R_x$ and $x\in Y_1(f)$, then $f$ is tangent
both to $R$ and to its selected coarse
$(\Delta,C_Rt)$-rectangle $\widetilde R$, provided
$C_R$ is sufficiently large.
\end{lemma}
\begin{proof}
\uses{lem:common-tangent,lem:cinematic-interval-growth}
Propagate the fine tangency bound across the longer
coarse interval, using the selected tangency scale
$\Delta$.
\end{proof}

Let $\mathcal G'(R)$ be the functions incident to $R$
after degree refinement. For a selected coarse rectangle
$\widetilde R$, let $F(\widetilde R)$ be the functions
incident to at least one of its retained fine rectangles.
The refinement makes the number of such incidences per
function comparable within each coarse rectangle.

\begin{lemma}[Many nondegenerate pairs; PYZ, Lemma 5.7]
\label{lem:good-function-pairs}
\lean{Kakeya.Cinematic.GoodPairCountingStatement}
For each retained fine rectangle $R$, at least
$\#\mathcal G'(R)^2/3$ ordered pairs $(f,g)$ satisfy
\[
 \delta^{2\eps}t\leq\|f-g\|_{C^2}\leq6t,
 \qquad
 \delta^\eps\Delta\leq\Delta(f,g)+\delta\leq2\Delta.
\]
\end{lemma}
\begin{proof}
\uses{lem:cinematic-two-ends}
The distance two-ends bound discards at most one third
of the pairs for being too close. The analogous
tangency selection discards at most another third.
The upper bounds follow from the selected common
distance and tangency bands.
\end{proof}

\begin{lemma}[One pair meets few fine rectangles; PYZ, Lemma 5.8]
\label{lem:pair-rectangle-count}
\lean{Kakeya.Cinematic.PairIncidenceCountingStatement}
A pair with
$\|f-g\|_{C^2}\geq\delta^{2\eps}t$ and
$\Delta(f,g)+\delta\geq\delta^\eps\Delta$
belongs to $\mathcal G'(R)^2$ for at most
$O(\delta^{-4\eps})$ retained fine rectangles $R$.
\end{lemma}
\begin{proof}
\uses{lem:tangency-sublevels,lem:rectangle-packing}
The common tangency is confined to at most two short
intervals. Cover these intervals by the permitted
rectangle locations and apply the incomparability
packing bound.
\end{proof}

\begin{lemma}[Nonconcentration in a coarse fibre; PYZ, Lemma 5.9]
\label{lem:coarse-fibre-nonconcentration}
For each retained coarse rectangle and every function-space
ball $B(g,r)$,
\[
 \#(F(\widetilde R)\cap B(g,r))\lesssim
       |\log\delta|(r/t)^\eps(t/\Delta)^{\eps^2}
                                      \#F(\widetilde R).
\]
\end{lemma}
\begin{proof}
\uses{lem:cinematic-two-ends}
Count incidences between functions in the ball and
fine rectangles inside $\widetilde R$. Apply the
pointwise two-ends bound at each assigned point, then
divide by the common refined degree. The tangency
selection and the degree pigeonhole account for the
last factor and the logarithm.
\end{proof}

\begin{remark}[A variable counting constant; PYZ, Remark 5.10]
\label{rem:pyz-variable-constant}
If the counting hypothesis on $F$ has a fixed factor
$A\geq1$ in place of its small-power factor, random
selection at rate $1/A$ gives the corresponding
$3/2$-moment bound with a factor $A^{1/2}$.
The same absolute normalization and threshold
dependencies must be retained when using this variant.
\end{remark}

